\subsection{WDM and moving-periodogram domains}
\label{sec:representations}

Both analyses fit the same tensor-product P-spline representation of
$\log S(t,f)$; they differ only in the observations supplied to the likelihood
(Fig.~\ref{fig:representations}).  We display the common domain using normalised
time $u=t/T\in[0,1]$ and normalised frequency
$f/f_{\mathrm{Nyq}}\in[0,1]$, where $f_{\mathrm{Nyq}}=1/(2\Delta t)$.
These axes describe the parent physical domain; the fitted WDM grid may exclude
boundary time bins and low- or high-frequency channels according to the
configured trims.  The WDM transform gives a critically sampled,
near-orthogonal tiling of the time--frequency plane with one real coefficient
per cell,
\begin{equation}
  w_{nm}\mid S_{nm}\sim\mathcal N(0,S_{nm}).
\end{equation}
This regular tiling uses the complete data record and localizes a disturbance in
both time and frequency.  For locally stationary coloured noise, treating the
coefficients as conditionally independent is an approximation whose adequacy
must be checked through the whitened residuals.

The Tang dynamic-Whittle construction instead considers a centred local DFT of
$2m+1$ real samples.  Its $2m+1$ outputs comprise one DC bin, $m$ positive-frequency
bins, and $m$ conjugate negative-frequency bins.  The construction evaluates
and retains only one positive-frequency bin at each window centre.  Its index is
$j(t)=1+[(t-1)\bmod m]$: thus the highlighted $f_3$ in
Fig.~\ref{fig:representations} is merely the example centre for which $j(t)=3$.
The selected bin cycles through $f_1,\ldots,f_m$, producing a diagonal block in the candidate
time--frequency grid.  Retained blocks begin $im$ window centres apart: after
the $m$ retained centres in one block, the intervening $(i-1)m$ centres are not
evaluated.  Because the local window must remain inside the record and only
positive-frequency bins are used, the retained Tang ordinates lie in the
interior $0<u<1$ and $0<f/f_{\mathrm{Nyq}}<1$.  Each retained squared Fourier
magnitude is modelled by the exponential Whittle likelihood
\begin{equation}
  I_t\mid S(t,f_t)\sim\operatorname{Exp}\!\left[
  \operatorname{mean}=S(t,f_t)\right].
\end{equation}
The cycling scheme covers the evolving spectrum without constructing a dense
periodogram at every time, while thinning limits the redundancy caused by
overlapping local windows.

Neither representation makes its finite-sample observations exactly
independent.  In the moving-periodogram analysis, approximate independence
follows from local stationarity, Fourier orthogonality, and thinning; in the WDM
analysis it follows from the near-diagonalization of a locally stationary
covariance by the wavelet transform.  We therefore regard agreement between the
two domains, together with representation-specific whitening checks, as a
robustness test rather than privileging either diagonal likelihood as exact.
The WDM localization is particularly useful in the presence of gaps, whereas a
gap can contaminate every Fourier ordinate whose $2m+1$ window intersects it.

\begin{figure}
  \centering
  \includegraphics[width=\linewidth]{representation_schematic.pdf}
  \caption{The two observation domains on
  $(u,f/f_{\mathrm{Nyq}})\in[0,1]^2$.  WDM produces a regular,
  critically sampled tiling with one real coefficient per
  $\Delta T\times\Delta F$ cell (left); configurable boundary trims are applied
  before fitting.  A real $2m+1$-sample local DFT contains
  DC, $m$ positive-frequency bins, and their $m$ conjugate partners.  Tang
  retains the single positive-frequency bin $f_{j(t)}$, cycling from $f_1$ to
  $f_m$ over $m$ successive window centres.  Retained blocks start $im$ centres
  apart, leaving $(i-1)m$ centres unevaluated (right).  Pale cells in the Tang
  panel denote candidate locations only; the implementation evaluates the
  highlighted ordinate directly rather than constructing a dense periodogram.
  The diagonal likelihoods
  in both domains are local-stationarity approximations, not claims of exact
  finite-sample independence.}
  \label{fig:representations}
\end{figure}
